% !TEX root = ../main.tex
\chapter{Spacetime}\label{ch:c13}
\begin{paracol}{2}

\begin{sources}
\begin{tabularx}{\linewidth}{@{}l l L@{}}
\toprule
\textbf{Source} & \textbf{Length} & \textbf{Content}\\
\midrule
\seg{13}{1}{Lecture 13: Spacetime} & 1:51:16 & Newtonian versus relativistic spacetime; the time orientation; the cones of the Lorentzian metric; the postulates for massive and massless particles; observers and their frames; the clock postulate and proper time; special relativity as the flat case and the twin paradox computed; a light clock; the spatial velocity measured by an observer; the Faraday tensor as seen by an observer; the role of the Lorentz transformations. Student questions on the role of $T$, time orientation, the frame bundle, the real world and charts, and the twin paradox.\\
\bottomrule
\end{tabularx}

\smallskip
\textbf{Transcript} \file{materials/transcripts/C13-V01.txt} (YouTube captions). There is no tutorial video for this lecture.
\end{sources}

\begin{reading}[Where this chapter sits in the other notes]
\emph{Theoretische Mechanik} (Schuller's classical mechanics notes): the dynamics and gravity of chapters 10 and 11 there are the Newtonian setting that this lecture replaces by the relativistic one; the Newtonian spacetime of chapter~\ref{ch:c09} is the comparison made all through the lecture [exact chapter titles not checked]. In this book the Lorentzian metric and signature are chapter~\ref{ch:c10}, the frame and its dual basis use chapters~\ref{ch:c03} and~\ref{ch:c05}, and the definition of relativistic spacetime was announced in chapter~\ref{ch:c01}. The next lecture (chapter~\ref{ch:c14}) introduces matter on this spacetime.
\end{reading}

% ==================================================================
\section{From Newtonian to relativistic spacetime}\label{sec:c13-newton}

Good morning and welcome to the International Winter School on Gravity and Light at the Johannes Kepler University in Linz \ts{13}{1}{0:00}. About 50 participants from all over the world join for the remainder of the school, the actual physical part. The first twelve lectures dealt with the mathematical foundations of general relativity; now the character of the course changes: the mathematical tools are employed to describe relativistic spacetime. One has to argue why one chooses which structure, and these arguments are rather vague, ``because it is about physics and physics is more complicated than mathematics'': mathematics is self-referential, and that is its power; it has nothing to do with the world and wants nothing to do with it; physics must relate its notions to the world and in the end make predictions, so more hand-waving and arguing is required, which is reduced to a minimum. The picture presented may need revision one day (maybe experiments already require it, and nobody knows how deep the revision will go); but to change a theory surgically one must understand the standard theory first. Only towards the end of the school will he say how the theory might be changed \ts{13}{1}{0:21}\ts{13}{1}{1:23}\ts{13}{1}{2:04}.

Recall the definition of \textbf{Newtonian spacetime} (\cref{def:c09-nst}): a set with a topology making it a four-dimensional manifold, with a smooth structure, a torsion-free connection and an absolute time, a smooth function $t$ on the manifold, with two compatibility conditions: $\mathrm dt$ vanishes nowhere (so there are time slices) and $\nabla\mathrm dt=0$ (time flows uniformly). Compare with the definition of relativistic spacetime given before lecture 1 (chapter~\ref{ch:c01}), which at the time was quite incomprehensible and now is not: it is a topological manifold of four dimensions, with a smooth (oriented) atlas and a torsion-free connection $\nabla$: so far the same. But where Newtonian spacetime had the absolute time, the relativistic one has two other structures: a \textbf{Lorentzian metric} $g$ (signature convention $(+,-,-,-)$, chapter~\ref{ch:c10}) and a \textbf{time orientation} $T$. What the absolute time function does in the Newtonian case is done here by the \emph{interplay} of these two. All other concepts in the long definition of lecture 1 are already clarified; the time orientation is the new one \ts{13}{1}{2:48}\ts{13}{1}{3:29}\ts{13}{1}{4:13}\ts{13}{1}{5:15}\ts{13}{1}{5:36}.

\begin{definition}[New concepts: time orientation, relativistic spacetime]\label{def:c13-timeorientation}
Let $(M,\mathcal O,\mathcal A,g)$ be a Lorentzian manifold (the connection plays no role yet), with an oriented atlas, which will be needed later. A \textbf{time orientation} is a vector field $T$ on $M$ that
\begin{enumerate}
  \item vanishes nowhere (``otherwise it is rather pointless'': a student points out that this condition had been forgotten on the board), and
  \item satisfies $g(T,T)>0$ at every point.
\end{enumerate}
A \textbf{relativistic spacetime} (from now on just \emph{spacetime}; ``Newtonian spacetime'' is used for the non-relativistic version) is a four-dimensional Lorentzian manifold with a smooth oriented atlas, a torsion-free connection and a time orientation.
\end{definition}
\ts{13}{1}{5:57}\ts{13}{1}{6:38}\ts{13}{1}{7:20}\ts{13}{1}{17:13}

\paragraph{The pictures: half spaces versus cones.} Back to Newtonian spacetime, in a picture with one dimension of space and one of time (``still too few, but I reduce the picture by two dimensions''). At a point the absolute time has a certain value and, being smooth, the same value on the surface of equal time. The gradient $\mathrm dt$ is a \emph{covector}: it is best understood not as an arrow but as a little plane attached at the point, an element of the cotangent space. It singles out the subset of the \emph{tangent} space at that point of the \textbf{future-directed} vectors $X$ with $\mathrm dt(X)>0$ (\cref{def:c09-fut}): a half space; the vectors lying in the plane are annihilated by $\mathrm dt$, those that point backwards have the opposite sign \ts{13}{1}{8:02}\ts{13}{1}{9:05}\ts{13}{1}{9:47}\ts{13}{1}{10:31}.

In relativistic spacetime there is no absolute time. What one has instead is the Lorentzian metric, which provides in each tangent space the double-cone structure seen in the tutorials: the vectors on the cone satisfy $g(X,X)=0$, those inside the double cone (up and down) satisfy $g(X,X)>0$, those outside $g(X,X)<0$. Future-directed particle world lines in Newtonian spacetime have future-directed tangents; here one wants to restrict to world lines whose tangent vectors lie inside the cone or possibly on its boundary, not running backwards and not lying outside. The motivation is the idea that \emph{no massive particle can overtake light}. (``High velocity'' is only in quotation marks for the moment: velocity is a vector, and without a metric it is not clear what a high velocity is.) One wants to distinguish \emph{one half} of the double cone, a half cone, as in Newtonian mechanics one has only a half space. The metric alone does not do this, it treats the two halves alike; this is what $T$ is for: a smooth vector field with $g(T,T)>0$ lies at each point either in the upper or in the lower cone, and by smoothness it cannot suddenly flip to the opposite cone, so it selects one half-cone in each tangent space \ts{13}{1}{11:14}\ts{13}{1}{11:34}\ts{13}{1}{12:19}\ts{13}{1}{13:01}\ts{13}{1}{13:41}\ts{13}{1}{14:43}\ts{13}{1}{15:04}.

So \emph{the combination of a Lorentzian metric and a time orientation gives a distribution of half-cones over all the tangent spaces of the spacetime}, as the absolute time gives a distribution of half-spaces in Newtonian mechanics. Cones live in the \emph{tangent spaces}, not in the manifold (they differ from point to point: at another point $Q$ there is another cone in $T_QM$) \ts{13}{1}{14:23}\ts{13}{1}{15:25}. In Newtonian spacetime one can always choose a stratified atlas in which the surfaces of equal time are horizontal; generically it is not possible to find an atlas in which all the cones stand upright, but cones standing upright or not is an artefact of the chart \ts{13}{1}{15:48}\ts{13}{1}{16:09}.
Schuller emphasises that he did \emph{not} argue that Newtonian physics has a metric on three-space and relativity one on four-space (one can introduce a metric in Newtonian physics too, in inertial systems): ``the first and foremost task of the Lorentzian metric is not to measure any length. The first and foremost task is to produce these cones'' \ts{13}{1}{16:29}\ts{13}{1}{16:50}.

\begin{coursenote}[Student question: what is the role of $T$?]
``I understand that from the cones we need a metric, but how does the $T$ come in?'' If one is given only a manifold with a Lorentzian metric, then in the tangent space at a point the vectors on the cone have $g(X,X)=0$, those in the upper cone and in the lower cone have $g(X,X)>0$, and those outside $g(X,X)<0$ (with the sign convention fixed). Requiring positive sign selects the interior of the upper \emph{and} of the lower cone. To define particles one wants to talk about future and past, particles running forward rather than backward, so one of the two cones has to be eliminated; the metric does not distinguish them, so one needs an extra choice: a vector field $T$ lying in one of them. That is the role of $T$ \ts{13}{1}{17:34}\ts{13}{1}{17:54}\ts{13}{1}{18:15}\ts{13}{1}{18:56}\ts{13}{1}{19:16}\ts{13}{1}{19:37}.
\end{coursenote}
\begin{coursenote}[Student question: how do we know which cone is the future one?]
Because $T$ says so. ``Relativists came up with no better idea than to say: choose it.'' Since $T$ is smooth, once chosen at one point the choice is more or less determined at the others (otherwise one just exchanges the notions of future and past). The arrow of time, whether time runs like this or like that, may have answers from statistical physics and other ideas; in general relativity this is not further resolved: ``we say, this is what it is, we have such a time orientation''. It is different from the oriented atlas: it is a choice of a half cone \ts{13}{1}{27:26}\ts{13}{1}{27:50}\ts{13}{1}{28:10}\ts{13}{1}{28:30}.
\end{coursenote}
\begin{coursenote}[Student question: can every manifold carry these structures?]
No. Every smooth manifold can carry a Riemannian metric, but not every one a Lorentzian metric: there are obstruction theorems which tell which can and which cannot, and the same holds for a time orientation. For instance the sphere $S^2$, which is orientable, cannot be made into a spacetime: one does not find a Lorentzian metric on it. So adding further structures constrains the underlying manifold, and one can enter the business of deriving theorems on what survives at the lower level \ts{13}{1}{28:52}\ts{13}{1}{29:14}\ts{13}{1}{29:34}\ts{13}{1}{29:54}\ts{13}{1}{30:15}.
\end{coursenote}
\begin{coursenote}[The inner product here is a pseudo inner product]
The Lorentzian metric is not an inner product in the sense of a positive-definite one, because it has the wrong signature (positive signature throughout would induce an inner product on each tangent space). It induces a \emph{pseudo inner product} on each tangent space: it shares all the properties but the positivity, and in certain directions an \emph{inverse triangle inequality} replaces the triangle inequality (``we will come to that''). ``If I say inner product I mean here Lorentzian, or pseudo, inner product'' \ts{13}{1}{24:08}\ts{13}{1}{24:30}.
\end{coursenote}

\begin{aside}[Cones as a generalisation of vector subspaces, and ``9/10 of a metric'']
A \emph{cone} is an object of linear algebra generalising a vector subspace: a subset where the sum of two vectors in it is in it, as for a subspace, and positive multiples of a vector in it are again in it, a stronger requirement than ``any multiple'' (which would give a subspace). A half space is a cone; the half-cone of the spacetime is a cone; both are even \emph{convex} cones. ``Relativity is: instead of having cones that are half spaces, you have cones that are proper cones'', and one may imagine the half-space opening up more and more until it looks like a narrower and narrower cone. Schuller: ``all the difficulties one might have at first with relativity is that somehow in your head you still have this absolute time function somewhere'' and cannot wrap one's head around the tangent spaces being divided not into two half spaces but into one cone and ``not that one cone''. ``That is all the problems with relativity, there are no others''; if one follows the definitions no such problem arises. He also says that \emph{spacetime is essentially this collection of cones}: ``in fact $9/10$ of a metric are determined by the cones'', only $1/10$ remains to be fixed (``so if I say spacetime is essentially a distribution of cones it is almost right, only $10\%$ error''). (Not said in class: a symmetric $(0,2)$ tensor in four dimensions has $10$ independent components; the cones fix the metric up to a positive function factor, i.e.\ all but one function.) \ts{13}{1}{30:57}\ts{13}{1}{31:19}\ts{13}{1}{31:40}\ts{13}{1}{32:00}\ts{13}{1}{32:42}\ts{13}{1}{33:03}\ts{13}{1}{33:24}\ts{13}{1}{34:06}
\end{aside}
Of course the whole idea is that these particles, ``if you think in velocities'', can never overtake light. Usually one says nothing runs faster than light, but faster refers to speed, and this has not been talked about yet: one could measure the length of the spacetime velocity $v$, but that is not the velocity an observer sees; speed means what an observer sees in his laboratory. This comes below \ts{13}{1}{34:28}\ts{13}{1}{34:48}\ts{13}{1}{35:09}.

\section{Postulates for the world lines of particles}\label{sec:c13-particles}

The definition is made to enable the following physical \emph{postulates}:

\begin{definition}[Postulate P1: massive particles]\label{def:c13-p1}
The world line $\gamma$ of a \emph{massive particle} satisfies, at every point $\gamma(\lambda)$ of the curve:
\begin{enumerate}
  \item $g\bigl(v_\gamma,v_\gamma\bigr)\big|_{\gamma(\lambda)}>0$, i.e.\ the tangent lies inside the (forward or backward) double cone;
  \item $g\bigl(T,v_\gamma\bigr)\big|_{\gamma(\lambda)}>0$, which selects the forward cone: it lies in the same cone as $T$.
\end{enumerate}
\end{definition}
\ts{13}{1}{19:58}\ts{13}{1}{20:46}\ts{13}{1}{21:28}
Picture: a smooth curve in $M$; at each point of the curve one checks whether its tangent vector (a directional derivative, an element of the tangent space, ``although I draw it like an arrow''), lies inside the cone that lives at that point; and in which cone: if $T$ selected the other cone, the inner product of the velocity with $T$ would be negative. One then checks the next point, with the cone \emph{at that point}: that is why it is important that spacetime consists of a half cone in each tangent space, the ones given by the metric and $T$ \ts{13}{1}{21:49}\ts{13}{1}{22:09}\ts{13}{1}{22:31}\ts{13}{1}{22:51}\ts{13}{1}{23:13}\ts{13}{1}{23:34}.

\begin{definition}[Postulate P2: massless particles]\label{def:c13-p2}
The world line $\gamma$ of a \emph{massless particle} satisfies at every point
\begin{enumerate}
  \item $g\bigl(v_\gamma,v_\gamma\bigr)\big|_{\gamma(\lambda)}=0$: the tangent lies on the boundary of the cone (in the captions the ``$=0$'' is said as ``the velocity is equal to zero'', a slip of speech);
  \item $g\bigl(T,v_\gamma\bigr)\big|_{\gamma(\lambda)}>0$, formally as for P1 (``positive pseudo inner product'').
\end{enumerate}
\end{definition}
\ts{13}{1}{24:30}\ts{13}{1}{24:52}\ts{13}{1}{25:33}\ts{13}{1}{25:53}
Note that in Newtonian mechanics massless particles could not be discussed: $F=ma$ requires a non-vanishing mass; it is in relativity that massless particles exist. Picture: at every point of a Lorentzian spacetime one draws a little half-cone; the world line of a massive particle runs inside the cones at each point, that of a massless one is tangent to the boundary; the cones may be positioned in all kinds of ways (``it could also start looking like this, why not?'') \ts{13}{1}{24:52}\ts{13}{1}{26:14}\ts{13}{1}{26:34}\ts{13}{1}{27:06}\ts{13}{1}{30:36}\ts{13}{1}{30:57}.

\section{Observers}\label{sec:c13-observers}

Observers are needed to talk about notions we are familiar with, namely what an observer sees \ts{13}{1}{35:09}\ts{13}{1}{35:29}.

\begin{definition}[New concepts: observer]\label{def:c13-observer}
An \textbf{observer} on a spacetime is a world line $\gamma$ with
\[
  g\bigl(v_\gamma,v_\gamma\bigr)>0\qquad\text{and}\qquad g\bigl(T,v_\gamma\bigr)>0,
\]
together with a choice of basis $e_0(\lambda),e_1(\lambda),e_2(\lambda),e_3(\lambda)$ of the tangent space $T_{\gamma(\lambda)}M$ (in $D$ dimensions, $D$ vectors), depending on the same parameter $\lambda$ as the curve, such that
\begin{enumerate}
  \item $g\bigl(e_a(\lambda),e_b(\lambda)\bigr)=\eta_{ab}$, with $\eta=\mathrm{diag}(1,-1,-1,-1)$ (``one could say they are also normal; more like a unit, but orthonormal for a Lorentzian metric''), and
  \item $e_0(\lambda)=v_\gamma(\lambda)$, the tangent to the curve.
\end{enumerate}
One writes an observer as $(\gamma,e)$. Precisely, an observer is a smooth curve in the frame bundle $LM$ over $M$.
\end{definition}
\ts{13}{1}{35:29}\ts{13}{1}{36:16}\ts{13}{1}{37:38}\ts{13}{1}{38:22}\ts{13}{1}{40:57}\ts{13}{1}{41:17}\ts{13}{1}{41:38}
So an observer is a world line like that of a massive particle: ``observers in relativity are always observers that could be realised as massive objects, because they have to carry with them a measurement apparatus''. Schuller does not like the traditional names \emph{time-like}, \emph{space-like}, \emph{light-like}: ``there is no time in the game so far'', the ``-like'' says it does not mean what it says, ``let us just not say it; later on you may reinstate that, like swear words if you have children'' \ts{13}{1}{36:36}\ts{13}{1}{36:58}\ts{13}{1}{37:18}.

The picture at a point $P=\gamma(\lambda_0)$: the tangent vector $e_0(\lambda_0)$ is fixed (``for the first basis element I do not have a choice'') and three further vectors $e_1,e_2,e_3$ pairwise $\eta$-orthogonal, with $g(e_0,e_0)=1$ and $g(e_\alpha,e_\alpha)=-1$ (``I only draw two more because my picture is at most three-dimensional''). At another point the same, but $e_0$ is fixed while $e_1,e_2$ are not: intuitively any rotation of them would do. An observer is more than a world line: he feels, observes quite different things depending on whether he is freely falling or freely falling \emph{and rotating}, and the rotation is captured only via the attached frame \ts{13}{1}{41:38}\ts{13}{1}{42:20}\ts{13}{1}{43:03}\ts{13}{1}{46:56}\ts{13}{1}{47:17}.

\begin{coursenote}[Student questions on observers]
\begin{itemize}
  \item \emph{Newtonian spacetime has the time as a function: did it live in three dimensions?} No, Newtonian spacetime was a four-dimensional manifold; the absolute time is a smooth function on it, and the level surfaces of constant absolute time are three-dimensional surfaces, the space at that time \ts{13}{1}{39:11}\ts{13}{1}{39:31}\ts{13}{1}{39:52}.
  \item \emph{Is it normalised to one by definition?} The tangent vector of an arbitrary curve is not normalised, but here it is, on purpose: one can always reparametrise, and it helps later with the definition of time. One could have written ``$g(v,v)=1>0$'' already in the first condition \ts{13}{1}{43:44}\ts{13}{1}{44:05}\ts{13}{1}{44:26}.
  \item \emph{Does the basis have to respect the signature?} It must: one could not find a basis in which $g$ is diagonal and which does not have the signature (it could be $5,-7,-3$, but then one rescales the vectors) \ts{13}{1}{44:51}.
  \item \emph{What is a frame bundle?} A manifold with, at each point $P$, as fibre not the tangent space but the space of all bases of $T_PM$: a point of the fibre is a choice of basis, a quadruplet of tangent vectors that constitutes a basis. A smooth curve in the frame bundle encodes that the basis does not jump but changes smoothly along the curve. The idea suffices here, the precise statement can be ignored \ts{13}{1}{45:11}\ts{13}{1}{45:32}\ts{13}{1}{45:52}\ts{13}{1}{46:14}.
  \item \emph{Would it be sufficient to say that the basis vectors, as functions of $\lambda$, are differentiable?} The question is what precisely that means, since each $\lambda$ goes to a different tangent space; one could say there are four smooth curves in the tangent bundle, which are additionally orthonormal in the above sense: probably a leaner way \ts{13}{1}{46:14}\ts{13}{1}{46:35}.
\end{itemize}
\end{coursenote}

\section{The clock postulate and proper time}\label{sec:c13-clock}

\begin{definition}[Postulate P3: the clock postulate]\label{def:c13-p3}
A clock carried by an observer $(\gamma,e)$ and switched on at $\gamma(\lambda_0)$ and off at $\gamma(\lambda_1)$ measures the time interval
\[
  \tau:=\int_{\lambda_0}^{\lambda_1}\mathrm d\lambda\,\sqrt{g\bigl(v_\gamma,v_\gamma\bigr)\big|_{\gamma(\lambda)}}.
\]
\end{definition}
\ts{13}{1}{47:37}\ts{13}{1}{48:21}\ts{13}{1}{49:02}
In relativists' jargon an \emph{event} is a point in spacetime (``pushing a button is localised in space and in time''): the clock is started at the event $\gamma(\lambda_0)$ and stopped at $\gamma(\lambda_1)$. The time measured is the length of this part of the world line in spacetime, as measured by the metric, which is nothing but the length functional of chapter~\ref{ch:c10}. This seems a rather arbitrary postulate, but one can support it with a light clock (below); and it defines which time this observer sees. Note: this is the \emph{first instance in the discussion of relativistic spacetime where the word ``time'' enters}: time is a derived notion, the time that a clock measures. It is called the \textbf{proper time} (``eigentime'') of the observer along that part of the world line, to distinguish it from Newtonian time; ``in this theory it is just time; time is what a clock measures'' \ts{13}{1}{49:23}\ts{13}{1}{49:44}\ts{13}{1}{50:06}\ts{13}{1}{50:27}\ts{13}{1}{50:48}\ts{13}{1}{51:09}\ts{13}{1}{58:50}\ts{13}{1}{59:11}.

Nothing in the definition suggests that another clock, starting at the same point and ending at the same endpoint along another path (that is also an observer), shows the same time: a priori no, and in fact \emph{it will not}, the time depends on the world line the clock takes. ``If people say the concept of time is changed in relativity, that is right, because there is no longer one absolute time.'' In Newtonian spacetime one could pick a single point $P$, without observer, second point or clock, and know the time associated with it (up to a rescaling of the absolute time, but once settled an event comes at a certain time), and the time difference of two events $P,Q$ is $t(P)-t(Q)$ for any path and any clock. Here it is entirely different: there is no time associated to an individual event, not even a time difference for two events; there is a time difference associated to a certain \emph{path} between two events that the clock takes \ts{13}{1}{51:29}\ts{13}{1}{51:50}\ts{13}{1}{52:11}\ts{13}{1}{52:31}\ts{13}{1}{52:52}\ts{13}{1}{53:34}\ts{13}{1}{54:15}\ts{13}{1}{54:36}.

\begin{coursenote}[Student question: is all this about the real world or about a chart? Is it chart independent?]
It is the real world: no component is used, no chart; the metric, the velocity, the cones are objective. For example a light bulb blinks for a millisecond at the point ``here and now''; light goes out spherically (``depending on your frame of reference very roughly this''), but the picture drawn is in the tangent space: the actual light travels along a null trajectory whose tangent at that point is what is drawn. A photon sent in some direction has that velocity, ``unimpressed by the choice of chart, like the velocity of the bird'': the position of the cone is an objective reality. The question is justified, because he always speaks of the manifold and these objects as ``the real world'' and of charts as what we make of it; of course the manifold is a concept too: ``we have no further access to the real world but to make concepts and then experiments that relate these concepts''. Even ``the bird has a velocity'' is already thinking, one cannot see velocity in the world. But the concept of velocity, which a child has, does not depend on the three numbers one is indoctrinated with in school: those are already in a chart. Later a student asks whether the world line is not subjective if others look at it differently: no; the history of the bird is the history of the bird, one may picture it in different charts; all these pictures are metaphors (``that is why I do not like these pictures'') \ts{13}{1}{54:56}\ts{13}{1}{55:18}\ts{13}{1}{55:39}\ts{13}{1}{56:20}\ts{13}{1}{56:40}\ts{13}{1}{57:01}\ts{13}{1}{57:21}\ts{13}{1}{57:42}\ts{13}{1}{58:03}\ts{13}{1}{1:32:49}\ts{13}{1}{1:33:11}\ts{13}{1}{1:33:31}.
\end{coursenote}

\section{Special relativity and the twin paradox}\label{sec:c13-sr}

\begin{example}[Flat spacetime: special relativity]\label{ex:c13-sr}
Take $M=\R^4$ with the standard topology and a smooth atlas containing the identity chart $x=\mathrm{id}_{\R^4}$, and the metric whose components in that chart are $g_{ij}=\eta_{ij}$. Then the Christoffel symbols of the Levi-Civita connection are zero everywhere in that chart, and the chart covers the whole manifold (an oriented atlas). Choose the time orientation, in the same chart, as $T=(1,0,0,0)$. Since the connection coefficients vanish, the Riemann curvature vanishes: the spacetime is \emph{flat}, and in this chart it looks particularly nice. This situation is called \textbf{special relativity} (``one should say it is the domain of special relativity''): a flat spacetime. Otherwise special relativity is not different from general relativity \ts{13}{1}{59:31}\ts{13}{1}{59:54}\ts{13}{1}{1:00:36}\ts{13}{1}{1:00:57}\ts{13}{1}{1:01:19}\ts{13}{1}{1:01:39}.
\end{example}

Consider two observers, given in this chart (``because I need to write it down'', the objects themselves are chart independent). The first, $\gamma$, has components $\gamma(\lambda)=(\lambda,0,0,0)$, $\lambda\in[0,1]$: its derivative is $(1,0,0,0)$, which has positive inner product with $T$, so it is an observer. The second, $\delta$, has a more adventurous route, for a parameter $\alpha\in(0,1)$ (both ends excluded),
\[
  \delta(\lambda)=\begin{cases}(\lambda,\;\alpha\lambda,\;0,\;0)&\lambda\le\tfrac12,\\[2pt](\lambda,\;\alpha(1-\lambda),\;0,\;0)&\lambda>\tfrac12.\end{cases}
\]
(A student corrects the second line: Schuller had written $\alpha\lambda$ in a different place and agrees on the board that at $\lambda=\tfrac12$ the spatial component is $\tfrac12\alpha$ and at $\lambda=1$ it is zero again: ``I think you are right, thank you''.) It goes away from $\gamma$ and comes back to meet him again. In this chart the cones everywhere stand upright (that is the choice of $\eta$ in this chart; in another chart the same metric may look different). For $\alpha\to1$ the curve goes closer and closer to the boundary of the cone, but never hits it; for $\alpha\ge1$ it is no longer an observer \ts{13}{1}{1:02:01}\ts{13}{1}{1:02:22}\ts{13}{1}{1:02:42}\ts{13}{1}{1:03:04}\ts{13}{1}{1:03:48}\ts{13}{1}{1:04:32}\ts{13}{1}{1:07:20}\ts{13}{1}{1:08:28}.
This is the \textbf{twin paradox}: two twins born together, one falls freely in space, the other falls freely in another direction because the mother kicked him in that direction (a thought experiment in space), then is kicked back and meets the other again; or rockets and so on \ts{13}{1}{1:04:54}\ts{13}{1}{1:05:14}. One calculates the time shown by each clock between the two meetings, with the clock postulate:
\begin{align*}
  \tau_\gamma&=\int_0^1\mathrm d\lambda\sqrt{\eta_{ij}\dot\gamma^i\dot\gamma^j}=\int_0^1\mathrm d\lambda\sqrt{1\cdot1\cdot1}=1,\\
  \tau_\delta&=\int_0^{1/2}\mathrm d\lambda\sqrt{1-\alpha^2}+\int_{1/2}^1\mathrm d\lambda\sqrt{1-(-\alpha)^2}=\sqrt{1-\alpha^2}
\end{align*}
(velocity of $\delta$ is $(1,\pm\alpha,0,0)$ with $\eta$ having $+1$ and $-1$: only the time and one space component survive, and the second half is twice the same term). The curve $\delta$ is not smooth at $\lambda=\tfrac12$, but the integral can be evaluated piecewise, or the corner can be smoothed out with the same result: ``it is just that this is the easiest to integrate'' \ts{13}{1}{1:05:34}\ts{13}{1}{1:05:58}\ts{13}{1}{1:06:18}\ts{13}{1}{1:06:39}\ts{13}{1}{1:08:48}\ts{13}{1}{1:09:08}\ts{13}{1}{1:09:29}\ts{13}{1}{1:09:50}\ts{13}{1}{1:10:10}.
So for $\gamma$ one year passes (``1'') and for the travelling twin only $\sqrt{1-\alpha^2}<1$ if $\alpha\neq0$: the closer $\alpha$ goes to $1$, the more the world line goes to the boundary of the cone (``we would almost say approaches the speed of light, but I have not shown that yet'') and the more the time can be driven to zero: ``if you only move away fast enough and come back fast enough you do not age much according to the clock that is being carried'' \ts{13}{1}{1:10:31}\ts{13}{1}{1:10:51}.

Schuller stresses that ``who talked about any relativity? This is the definition of time shown by a clock'': at this point one may get suspicious of the whole framework, and indeed, taking postulate P3 seriously, one had better come up with a realistic clock design that supports it \ts{13}{1}{1:11:13}\ts{13}{1}{1:11:33}\ts{13}{1}{1:11:56}.

\begin{example}[A light clock]\label{ex:c13-lightclock}
Two little mirrors with a photon running between them, reflecting back and forth. Its history in spacetime: at each point there is the cone; the mirrors' world lines run inside the cones and the light on their boundary. It is a clock, it makes ``tick tock'', and the closer the mirrors, the more precise (``a fantastic clock, it works relativistically; light is a relativistic matter theory, we will come to that''). Take a second such clock which runs along the path of the travelling twin: count the ticks along the straight path of $\gamma$ ($12$ ticks, in the drawing) and along the other path ($8$ ticks): shorter. This is rough (what the real distance of the mirrors is in the moving system needs the terminology of the later lectures), but at this point the picture is more helpful than the abstract result alone \ts{13}{1}{1:12:16}\ts{13}{1}{1:12:41}\ts{13}{1}{1:13:03}\ts{13}{1}{1:13:23}\ts{13}{1}{1:13:43}\ts{13}{1}{1:14:04}\ts{13}{1}{1:14:25}\ts{13}{1}{1:14:46}\ts{13}{1}{1:15:07}.
\end{example}
\begin{coursenote}[Student questions on the twin paradox]
\begin{itemize}
  \item \emph{Did the example use the chosen basis of the observer?} No: the clock is only the length along the world line, independent of how the frame rotates or not (``we will get the notion of that''); the chart was used only to write the lines down \ts{13}{1}{1:15:27}\ts{13}{1}{1:15:48}\ts{13}{1}{1:16:09}.
  \item \emph{Each twin sees the other one move away; why does one age more?} ``Because of the clock postulate'' (``that was a mean answer''). The point: the one has inertial motion and the other has not: this one will realise at some point that he is being accelerated (``give both an iPhone and at some point the screen will change for this guy''), but that explains only why they are different, not the effect. It is sometimes claimed that all the effect is in the acceleration period: ``real nonsense, for many reasons''; it is the total length of the path that counts \ts{13}{1}{1:16:30}\ts{13}{1}{1:16:51}\ts{13}{1}{1:17:12}\ts{13}{1}{1:17:33}.
  \item \emph{Who has aged more if the twins go different ways and do not meet?} A totally ill-defined question: they have to reunite to compare their clocks. In flat spacetime each one sees as equal to his time the points on one family of hyperplanes, the other another family, and each asks ``now that I am this old, how old is he?''. So ``the whole effect appears here'' is a statement ``from whose point of view?''; clocks have to be synchronised and compared at the same point, because otherwise one does not even know when to compare \ts{13}{1}{1:17:33}\ts{13}{1}{1:17:53}\ts{13}{1}{1:18:14}\ts{13}{1}{1:18:34}\ts{13}{1}{1:18:55}.
  \item \emph{What if they could see each other during the movement?} Interesting questions, but every discussion now would rely on the spacetime being flat, where one can say which points are of equal time; in curved spacetime that is not even clear (there are notions like radar distance), and he does not want to confuse the audience with flat-space concepts \ts{13}{1}{1:19:16}\ts{13}{1}{1:19:36}\ts{13}{1}{1:19:56}.
\end{itemize}
\end{coursenote}

\section{The spatial velocity measured by an observer}\label{sec:c13-velocity}

\begin{definition}[Postulate P4: the three-velocity seen by an observer]\label{def:c13-p4}
Let $(\gamma,e)$ be an observer and $\delta$ a world line of a massive particle, parametrised such that $g(v_\delta,v_\delta)=1$ everywhere (``I parameterise world lines according to the time that a clock carried by the particle would show; this saves a lot of normalising factors, and one can always reparametrise''). Suppose they \emph{meet} somewhere in spacetime: there are parameters $\lambda_1$ for $\gamma$ and $\tau_2$ for $\delta$ with $\gamma(\lambda_1)=\delta(\tau_2)$ (meeting is a spacetime concept: ``I was here, you were here, but not at the same time''). The \textbf{spatial velocity} (three-velocity) of the particle measured by the observer is the vector
\[
  \mathsf v:=\sum_{\alpha=1}^{3}\varepsilon^\alpha\bigl(v_\delta(\tau_2)\bigr)\,e_\alpha(\lambda_1)
\]
where $(\varepsilon^0,\varepsilon^1,\varepsilon^2,\varepsilon^3)$ is the unique dual basis of $(e_0,e_1,e_2,e_3)$ at that point.
\end{definition}
\ts{13}{1}{1:20:20}\ts{13}{1}{1:21:02}\ts{13}{1}{1:21:47}\ts{13}{1}{1:22:07}\ts{13}{1}{1:23:35}\ts{13}{1}{1:24:21}\ts{13}{1}{1:24:46}\ts{13}{1}{1:25:09}
The velocity $v_\delta$ of the world line is the space-time velocity, sometimes called the \emph{four-velocity}, an objective tangent vector in the real world; but that is not the velocity spoken of in real life, which is a three-velocity. Only the three basis vectors $e_1,e_2,e_3$ enter the expansion, and $e_0$ is not part of it, so the resulting vector lies in the span of $e_1,e_2,e_3$ in the tangent space. The $\varepsilon^\alpha$ must be taken from the \emph{full} dual basis of four vectors and then only three of them are used: $\varepsilon^1,\varepsilon^2,\varepsilon^3$ is \emph{not} ``the unique dual basis to $e_1,e_2,e_3$'', because three vectors in a four-dimensional space are not a basis and have no unique dual basis \ts{13}{1}{1:25:31}\ts{13}{1}{1:26:12}\ts{13}{1}{1:26:34}\ts{13}{1}{1:26:55}.
It is observer dependent, ``for another observer it would look different, oh yes, and that is the whole point'': another observer with another frame extracts from the same objective four-velocity another three-velocity: ``that is the whole secret to different observers seeing different things''; it is not even very different from classical mechanics, where observers see different velocities too. The new thing is that, because of the cones, different observers can never see a massive particle with a velocity bigger than the speed of light. The spatial metric that the observer uses, in the tangent space and not on an extended submanifold, is $-g(e_\alpha,e_\beta)=\delta_{\alpha\beta}$ (a minus sign before it makes it friendlier); measuring $\mathsf v$ with it, Schuller says that this length ``will never exceed one'' and hence there is an upper limit to the spatial velocities one can see \ts{13}{1}{1:27:16}\ts{13}{1}{1:27:38}\ts{13}{1}{1:28:19}\ts{13}{1}{1:28:40}\ts{13}{1}{1:29:01}\ts{13}{1}{1:29:21}\ts{13}{1}{1:29:41}\ts{13}{1}{1:30:02}.
\begin{coursenote}[A caveat on the bound]
(Not said in class, flagged as a possible slip.) With $\mathsf v$ exactly as written, $v_\delta=\varepsilon^0(v_\delta)\,e_0+\mathsf v$ and $g(v_\delta,v_\delta)=1$ give $|\mathsf v|^2=\varepsilon^0(v_\delta)^2-1$, which is not bounded by one; the bound by one holds for the usual velocity $\mathsf v/\varepsilon^0(v_\delta)$, whose squared length is $1-1/\varepsilon^0(v_\delta)^2<1$. The lecture presumably means this normalised vector [?].
\end{coursenote}
The observer is a kind of very localised chart, only at one point: a frame in the tangent space (one might extend it to a chart on the manifold, but that is not even interesting). The concept is objective, derived from an objective one, namely the tangent to the world line; \emph{the three-velocity is a derived concept with respect to an observer}, and the objective thing is the world line and its tangent, the rest is extraction of information. ``One of the issues is getting rid of the notion that you somehow start objectively in three-space and then construct the four notions out of it. It is the wrong way around: you first have the notions in spacetime, they are objective, and then you construct the observer-dependent notions as derived objects'' \ts{13}{1}{1:30:22}\ts{13}{1}{1:30:42}\ts{13}{1}{1:31:03}\ts{13}{1}{1:31:23}\ts{13}{1}{1:31:44}.

\begin{coursenote}[Student questions on the spatial velocity]
\begin{itemize}
  \item \emph{Is the small $\mathsf v$ only the absolute value of the arrow?} No: it is a vector, the expansion in $e_1,e_2,e_3$: the $\varepsilon^\alpha$ are covectors acting on the four-vector; $\mathsf v$ is even a vector of spacetime, which just happens to lie in the span of $e_1,e_2,e_3$ \ts{13}{1}{1:31:44}\ts{13}{1}{1:32:04}.
  \item \emph{Why a postulate, after defining proper time? Can it be connected to the velocity of classical mechanics?} It looks funny to postulate, but what one would need is, for an observer at a point, all the points that are at the same time from his point of view, then one could consider curves through them and define these velocities. But the observer is just the tangent vectors in this and that direction and says nothing about what happens a little bit away in the manifold; in curved spacetime that is an issue of definition: how far is somebody away? One may send light that is reflected and comes back and use the \emph{radar time} (half of the time from here to there) as a measure of distance; but if some vicious entity (a black hole) removes the point where the light is reflected, the ray never comes back: is the other guy infinitely far away? So these ideas of distance are a matter of definition and one must think carefully about what is meant. ``Locally it will probably work'', but defining first space and then curves in it is less clean than this: where I am I know that I am there, the tangent space is there, I have the frame \ts{13}{1}{1:27:58}\ts{13}{1}{1:33:51}\ts{13}{1}{1:34:12}\ts{13}{1}{1:34:32}\ts{13}{1}{1:34:52}\ts{13}{1}{1:35:14}\ts{13}{1}{1:35:34}\ts{13}{1}{1:35:54}\ts{13}{1}{1:36:15}\ts{13}{1}{1:36:35}.
  \item \emph{This assumes the same point; what if the curvature is so large that one cannot do the experiment?} If the curvature gets very high it gets more and more difficult; this is taken up later \ts{13}{1}{1:36:56}\ts{13}{1}{1:37:17}.
\end{itemize}
\end{coursenote}

\section{The Faraday tensor as seen by an observer}\label{sec:c13-faraday}

An observer extracts quantities measurable in his laboratory from objective spacetime quantities always in this way. The prime example is the \textbf{Faraday tensor}, the $(0,2)$ tensor $F$ of electromagnetism. In a chart one has seen the antisymmetric array with zeros on the diagonal, the electric components $E_1,E_2,E_3$ in the first row (and minus these in the first column) and the magnetic components $B_1,B_2,B_3$ in the spatial block: ``the tensor that combines the $E$ and the $B$ field''. But this is already $F(e_a,e_b)$ \emph{in the frame of a particular observer}: ``no wonder that different observers see different $E$ and $B$ fields''. One extracts
\[
  E_\alpha:=F(e_0,e_\alpha),\qquad B_\gamma:=\tfrac12\,\epsilon_{\alpha\beta\gamma}\,F(e_\alpha,e_\beta)
\]
[up to signs and conventions; the formulas are only described aloud: ``you only look at the $e_0e_\alpha$ component and call it $E_\alpha$, you only look at the $e_\alpha e_\beta$ component, take an $\epsilon_{\alpha\beta\gamma}$ and call that $B_\gamma$''], with $\epsilon_{123}=+1$ totally antisymmetric. The objective spacetime quantity is the Faraday $(0,2)$ tensor, unimpressed by charts or observers; it is what the Maxwell equations say it is. In the lab one measures $E,B$ with respect to the lab frame: ``your lab is an observer; your lab may be rotating, for instance; then you measure different $E$ and $B$ fields''. So the observer is the \emph{intertwiner} between the objective quantities of the real world, described by abstract tensors, and the quantities measured in the lab with respect to the observer's frame. This fully explains the transformation law of $E$ and $B$ when one changes observer; in the $\epsilon$ one sees some funny behaviour of $B$ fields. A warning: $\epsilon_{\alpha\beta\gamma}$, a pure number, appears; each $E_\alpha$ is a function, but something that is a spacetime tensor may in general be a \emph{hypersurface density}: ``as soon as you start relating things to three-spaces, they might change their density weight'' (densities were mentioned before; this is for the interpretation) \ts{13}{1}{1:37:38}\ts{13}{1}{1:38:02}\ts{13}{1}{1:38:23}\ts{13}{1}{1:38:45}\ts{13}{1}{1:39:27}\ts{13}{1}{1:39:49}\ts{13}{1}{1:40:30}\ts{13}{1}{1:40:51}\ts{13}{1}{1:41:11}\ts{13}{1}{1:41:52}\ts{13}{1}{1:42:12}.

\section{The role of the Lorentz transformations}\label{sec:c13-lorentz}

The last point. The Lorentz transformations emerge as follows \ts{13}{1}{1:42:54}\ts{13}{1}{1:43:20}.

\begin{theorem}[Observer frames at a common point are related by Lorentz transformations]\label{thm:c13-lorentz}
Let $(\gamma,e)$ and $(\tilde\gamma,\tilde e)$ be two observers that meet, $\gamma(\tau_1)=\tilde\gamma(\tau_2)=:P$. Then there is a matrix $\Lambda\in GL(4)$ with $\tilde e_a=\Lambda^m{}_a\,e_m$ and
\[
  \eta_{ab}=\Lambda^m{}_a\,\Lambda^n{}_b\,\eta_{mn},
\]
i.e.\ $\Lambda$ is a \textbf{Lorentz transformation}.
\end{theorem}
\ts{13}{1}{1:43:47}\ts{13}{1}{1:44:08}
\begin{proof}
For simplicity choose the parameters such that the meeting is at $\tau=0$ for both. The frames $(e_0,\dots,e_3)$ and $(\tilde e_0,\dots,\tilde e_3)$ at $P$ are both bases of the \emph{same} tangent space $T_PM$ (that is the definition, ``totally clear''). Hence one can express the twiddle basis in terms of the untwiddled one: there is a change of basis $\Lambda$, a matrix in $GL(4)$ (a general linear, invertible map, otherwise it could not send a basis to a basis), $\tilde e_a=\Lambda^m{}_a e_m$. By the definition of an observer, $g(\tilde e_a,\tilde e_b)=\eta_{ab}$. Insert the expansion and use linearity of $g$ (even $C^\infty$-linearity, though plain linearity is enough at a point):
\[
  \eta_{ab}=g\bigl(\Lambda^m{}_ae_m,\Lambda^n{}_be_n\bigr)=\Lambda^m{}_a\Lambda^n{}_b\,g(e_m,e_n)=\Lambda^m{}_a\Lambda^n{}_b\,\eta_{mn},
\]
again because $(\gamma,e)$ is an observer. The matrices $\Lambda$ satisfying this condition are what is usually called the Lorentz transformations.
\end{proof}
\ts{13}{1}{1:44:30}\ts{13}{1}{1:44:50}\ts{13}{1}{1:45:35}\ts{13}{1}{1:45:57}\ts{13}{1}{1:46:17}\ts{13}{1}{1:46:37}\ts{13}{1}{1:46:58}
(In the lecture the proof is given in this order, with the computation first started in components of $g$ and then redone as above.)

\paragraph{The role of Lorentz transformations.} The first thing: they \emph{relate the frames of any two observers at the same point}. It is perfectly meaningless to talk about the frame of one observer and the frame of another and say that one comes from the other by a Lorentz transformation if they live in different tangent spaces. Lorentz transformations do \emph{not} act on spacetime: they act on \emph{one tangent space} of spacetime and relate observers who meet there; or they act on another tangent space and relate frames of observers who meet there. This has nothing to do with acting on spacetime points: ``the $\Lambda^\mu{}_\nu x^\mu$ is the new spacetime point $\tilde x$: they all say, utter nonsense, this just is not done'' \ts{13}{1}{1:47:19}\ts{13}{1}{1:47:42}\ts{13}{1}{1:48:03}\ts{13}{1}{1:48:25}\ts{13}{1}{1:48:46}.
One can do so in flat spacetime and say that all transformations one wants to look at are those; but then one greatly restricts to linear transformations (one never has to look at $\Gamma$'s if the spacetime was flat and one chose the right chart) and one cannot go to polar coordinates, and ``what has that to do with relativity? Absolutely nothing''. So Lorentz transformations \emph{on spacetime} are a terrible conceptual idea. One does that in special relativity, and then says that general relativity now allows not only Lorentz transformations but any transformation, or that ``the Lorentz group is replaced by the diffeomorphism group in general relativity'': this was already allowed in special relativity. Both in special and in general relativity, Lorentz transformations relate observer frames at the same point, whether the spacetime is flat or curved; general diffeomorphisms can be chosen to go to different charts, and there is no difference. ``Hence it is sometimes so dangerous to start with special relativity and concepts of the Minkowski vector space as the space of events, because it clouds your thinking: it is just an abuse of the flat-space structure'' \ts{13}{1}{1:49:06}\ts{13}{1}{1:49:27}\ts{13}{1}{1:49:50}\ts{13}{1}{1:50:10}\ts{13}{1}{1:50:31}\ts{13}{1}{1:50:52}.

\section*{Key concepts}
\begin{keyconcepts}
\textbf{Retrieval questions.}
\begin{enumerate}
  \item Compare Newtonian and relativistic spacetime: which structures replace the absolute time? What does each of the Lorentzian metric and the time orientation do, and why is neither enough alone?
  \item Define a time orientation. Why does the sphere $S^2$ not carry a spacetime structure?
  \item State the postulates P1 and P2 for the world lines of massive and massless particles. Where does each world line lie relative to the cones?
  \item Define an observer, including its frame. Why is it not just a world line?
  \item State the clock postulate. Why is ``time'' a derived notion? Compute the proper times of the twin-paradox example.
  \item How does an observer extract a spatial velocity from the objective four-velocity? Why is it observer dependent?
  \item How does an observer extract $E$ and $B$ from the Faraday tensor?
  \item Prove that two observers' frames at a common point are related by a Lorentz transformation. Why do Lorentz transformations not act on spacetime points?
\end{enumerate}
\textbf{Concepts.}
\begin{itemize}
  \item Relativistic spacetime, time orientation $T$, cones in tangent spaces (\cref{def:c13-timeorientation}).
  \item Postulates P1, P2 for particle world lines (\cref{def:c13-p1}, \cref{def:c13-p2}).
  \item Observer $(\gamma,e)$, frame bundle (\cref{def:c13-observer}).
  \item Clock postulate, proper time $\tau=\int\sqrt{g(v,v)}\,\mathrm d\lambda$; twin paradox (\cref{def:c13-p3}, \cref{ex:c13-sr}).
  \item Spatial velocity, $E$ and $B$ as extracted from $F$ (\cref{def:c13-p4}).
  \item Lorentz transformations as changes of frame at one point (\cref{thm:c13-lorentz}).
\end{itemize}
\end{keyconcepts}

\flushside
\end{paracol}
